%%
%% Chapter 1: Introduction
%% Independent file, included by ../main.tex
%% This draft is centered on the GT-conditioned CVAE algorithm.
%%

\section{Introduction}
\label{sec:intro}

Cone-beam computed tomography (CBCT) is a cornerstone of image-guided
procedures in radiation therapy and interventional radiology, offering
volumetric imaging with compact gantry geometries. Clinical workflows,
however, place hard limits on the number of projection views that can be
acquired: fewer views reduce radiation dose and shorten scan times, but
dramatically undersample the Fourier domain of the object. Sparse-view CBCT
reconstruction -- recovery of a high-quality volume from a handful of
projections, typically far below the hundreds required by analytical
reconstruction -- is therefore an ill-posed inverse problem whose solution
requires strong prior knowledge of anatomy.

Deep learning has become the dominant tool for this task. Neural
reconstruction methods range from direct 2D$\to$3D back-projection networks
to score-based and diffusion models, and recent work has shown that
explicit scene representations such as 3D Gaussian splatting (3DGS)
\cite{kerbl2023gaussiansplats} provide an outstanding inductive bias for
tomographic volumes: compact, differentiable, and easily supervised by
projection consistency. The DIF-Gaussian framework \cite{lin2024difgaussian}
demonstrated that a splatting-based reconstructor can produce accurate
volumes from as few as six cone-beam projections.

Despite these advances, the vast majority of sparse-view CBCT models are
\emph{deterministic}: they map a fixed set of views to a single volume and
output no measure of confidence. This is unsatisfactory in a clinical
setting, where reconstruction noise and ambiguity are spatially
inhomogeneous and task-relevant, and where downstream planning decisions
(e.g., dose painting or target delineation) should ideally consume
calibrated uncertainty, not a point estimate.

A natural remedy is a \emph{conditional variational} formulation. During
training, privileged information that will not be available at deployment --
the dense-view set, or even the ground-truth volume itself -- can be
encoded into a variational posterior $q(z\mid C,Y)$ to concentrate the
latent code on target-specific content and tighten the evidence lower
bound. At deployment this posterior is discarded, and the reconstruction is
generated by sampling a sparse-view prior $p(z\mid C)$. If the training and
deployment information flows are carefully separated, such a design yields
both accuracy and per-point uncertainty from the very same model. If they
are not separated, a subtle and common failure appears: when the training
decoder can read dense-view features directly, the latent becomes redundant
and the variational posterior \emph{collapses} onto an uninformative prior,
silently degrading the deployment path.

In this paper we design, build, and rigorously evaluate a
\emph{GT-conditioned CVAE} on top of the Gaussian-splatting CBCT
reconstructor. Our contributions are:

\begin{enumerate}
\item \textbf{A conditional variational formulation for splatting-based
CBCT reconstruction.} We inject a tractable 64-dimensional global latent
variable into the Gaussian-splatting decoder through feature-wise linear
modulation (FiLM), with a neural-process prior $p(z\mid C)$ computed from
the sparse context views $C$ alone.

\item \textbf{The GT-conditioned CVAE.} During training only, a
lightweight 3D CNN ($1\to16\to16\to32$ channels) encodes the ground-truth
volume into a 32-dimensional descriptor that conditions the variational
posterior $q(z\mid C,Y)$; the objective simultaneously supervises the
deployment-side prior. At inference the GT encoder and posterior are
discarded and $z\sim p(z\mid C)$; \emph{no} ground truth or dense views are
required, and per-point predictive variance comes for free from $N{=}8$
Monte-Carlo sampling.

\item \textbf{Empirical analysis of the decisive design choices.} The
GT-conditioned posterior matches the deterministic backbone within
$0.04$\,dB while adding per-point variance, and ablations isolate two
factors that dominate performance: the variational posterior's access
to the ground truth and the token initialization of the fusion
transformer ($\sim\!2$\,dB). The decoder is routed through the sparse
context features so the latent remains the sole channel of dense
information.

\item \textbf{Sparse-view results with calibrated uncertainty.} On the
LUNA16 CBCT benchmark with six input views, the GT-conditioned CVAE
reaches $29.50$\,dB PSNR / $0.868$\,SSIM, within $0.04$\,dB of the
deterministic backbone ($29.54$\,dB / $0.870$\,SSIM), while
additionally producing deployment-side variance estimates for every
reconstructed point.
\end{enumerate}

The remainder of this paper is organized as follows. Section~\ref{sec:rel}
reviews related work on sparse-view CT reconstruction, Gaussian splatting,
and conditional latent-variable models. Section~\ref{sec:method} presents
the proposed probabilistic formulation in detail.
Section~\ref{sec:exp} describes the experimental protocol, results, and
ablations. Section~\ref{sec:conclusion} concludes and outlines future
directions.